% SGA TRANSLATION PROJECT
% Volume: SGA 4
% Expose: IX
% Range translated: introduction and Sections 1--2, source lines 53--1135 of 09/09.tex.
% Status: draft translation, compile-checked, not proofed line-by-line against scan.

\begin{center}
{\Large \textbf{SGA 4, Expos\'e IX}}\\[0.35em]
{\large \textbf{Constructible Sheaves; Cohomology of an Algebraic Curve}}\\[0.25em]
{\normalsize M. Artin}
\end{center}

\bigskip

\section*{0. Introduction}

Starting with the present expos\'e, in contrast with the preceding expos\'es and with the classical cohomological theory of topological spaces, we are forced, for the validity of the essential statements, to restrict to \emph{torsion} sheaves, and often more precisely to torsion sheaves prime to the residual characteristics.

Let us note that in fact a ``reasonable'' theory of cohomology with integral, or even real, coefficients for varieties over an algebraically closed field of characteristic \(p\ne 0\), analogous to the classical theory for \(k=\Cc\), does not exist, as Serre observed. More precisely, there is no contravariant functor \(H^1\), defined for instance on the category of smooth projective schemes over an algebraically closed field \(K\) of characteristic \(p>0\), with values in finite-dimensional vector spaces over \(\Rr\), or over a subfield of \(\Rr\), commuting with products, \ie such that
\[
H^1(X\times Y)\xleftarrow{\sim}H^1(X)\times H^1(Y),
\]
and such that \(\dim H^1(X)=2\) when \(X\) is a connected curve of genus 1. Indeed, it would follow that, if \(X\) is an abelian variety over \(k\), the map \(u\mapsto u^1\) from \(\End(X)\) to \(\End(H^1(X))\) is additive, hence is a representation of the opposite ring of \(\End(X)\). But in characteristic \(p\), there exist elliptic curves \(X\) whose endomorphism ring is a maximal order \(A\) in a quaternion algebra over \(\Zz\) ([3], p. 198), and such an algebra plainly has no representation on a vector space of dimension 2 over \(\Rr\), since \(A\otimes_{\Zz}\Rr\) is a division ring, namely Hamilton's quaternions.

\section*{1. The formalism of torsion sheaves}

Let \(T\) be a topos and \(F\) an abelian sheaf on \(T\). We can multiply by \(n\) in \(F\), for \(n\in\Zz\); this multiplication is induced by the analogous operation in \((\mathbf{Ab})\). Denote by \({}_nF\) the kernel of this multiplication. Thus \({}_nF(X)\), for \(X\in\Ob T\), is the group of sections of \(F(X)\) whose order divides \(n\). Let \(\mathcal P\) denote the set of all prime numbers.

\begin{statement}{Definition 1.1}
Let \(p\) be a set of prime numbers. We say that \(F\) is a sheaf of \(p\)-torsion, or a \(p\)-sheaf, if the canonical morphism
\[
\varinjlim_n {}_nF\longrightarrow F
\]
is bijective, where \(n\) runs through the positive integers such that \(\Ass(n)\subset p\), \ie such that all prime divisors of \(n\) belong to \(p\). If \(p=\mathcal P\) is the set of all prime numbers, we simply say that \(F\) is torsion.
\end{statement}

\begin{statement}{Proposition 1.2}
\begin{enumerate}[label=(\roman*)]
\item The sheaf \(F\) is of \(p\)-torsion if and only if \(F\) is the sheaf associated with a presheaf with values in abelian groups of \(p\)-torsion. One may take
\[
P=\varinjlim_{\Ass(n)\subset p}^{\mathrm{presheaves}}{}_nF.
\]
\item If \(F\) is of \(p\)-torsion and \(X\in\Ob T\) is quasi-compact, then \(F(X)\) is a group of \(p\)-torsion.
\item If \(T\) is locally of finite type (VI, 1.1), then \(F\) is of \(p\)-torsion if and only if \(F(X)\) is a group of \(p\)-torsion for each quasi-compact \(X\in\Ob T\). In this case one has
\[
\varinjlim_{\Ass(n)\subset p}H^q(T/X;{}_nF)\xrightarrow{\sim}H^q(T/X;F)
\]
for every quasi-compact \(X\) and every \(q\).
\item If \(u:T\to T'\) is a morphism of topoi and if \(F'\) is of \(p\)-torsion on \(T'\), then the inverse image \(u^*F'\) is a sheaf of \(p\)-torsion on \(T\).
\item If \(u:T\to T'\) is a morphism of topoi, with \(T\) and \(T'\) locally of finite type and \(u\) quasi-compact, and if \(F\) is a sheaf of \(p\)-torsion on \(T\), then the sheaves \(R^qu_*F\) are of \(p\)-torsion on \(T'\), for all \(q\).
\end{enumerate}
\end{statement}

\begin{prooftext}
For (i), if a sheaf \(F\) is of \(p\)-torsion, then \(F\) is plainly the sheaf associated with the presheaf \(P\) whose group of sections \(P(X)\) is the subgroup of \(p\)-torsion elements of \(F(X)\). Conversely, let \(F=aP\) be the sheaf associated with \(P\), where \(P\) takes values in abelian groups of \(p\)-torsion. From the exact sequence
\[
0\longrightarrow {}_nP\longrightarrow P\xrightarrow{n}P
\]
one obtains an exact sequence
\[
0\longrightarrow a({}_nP)\longrightarrow F\xrightarrow{n}F,
\]
hence \(a({}_nP)={}_nF\). Since the functor \(a\) commutes with inductive limits, the isomorphism
\[
\varinjlim_{\Ass(n)\subset p}^{\mathrm{presheaves}}{}_nP\xrightarrow{\sim}P
\]
gives
\[
\varinjlim_{\Ass(n)\subset p}{}_nF\xrightarrow{\sim}F.
\]

For (ii), put \(P=\varinjlim^{\mathrm{presheaves}}{}_nF\). Then \(P\subset F\), so \(P\) is a separated presheaf, and
\[
aP(X)=\varinjlim \ker\Bigl(\prod_iP(X_i)\rightrightarrows\cdots\Bigr),
\]
where, as usual, the limit is taken over the covering families \(\{X_i\to X\}\). Since \(X\) is quasi-compact, it is enough to use finite covering families. For such families, \(\prod_iP(X_i)\) is of \(p\)-torsion, and hence so is its kernel. Thus \(aP(X)=F(X)\) is of \(p\)-torsion.

For (iii), to check that \(\varinjlim {}_nF\to F\) is bijective, it is enough to check values on each \(X\) in a family of generators. Thus the first assertion reduces to (ii). The displayed isomorphism follows from the same passage to inductive limits in cohomology for locally finite type topoi. Assertion (iv) follows from (i), since \(u^*(aP')=a(u^\bullet P')\). Assertion (v) is an easy consequence of (iii).
\end{prooftext}

We shall also need a non-abelian variant. There are technical difficulties in the definition, and we content ourselves with giving the definition for the \'etale topology of a scheme.

\begin{statement}{Definition 1.3}
Let \(p\) be a set of prime numbers. A group \(G\) is called an \emph{ind-\(p\)-group} if every finite subset of \(G\) generates a finite subgroup of order \(n\), with \(\Ass(n)\subset p\). Equivalently, the finite subgroups \(G_\alpha\) of order \(n_\alpha\), with \(\Ass(n_\alpha)\subset p\), form a filtered system and \(G=\varinjlim G_\alpha\). If \(p=\mathcal P\), one says \emph{ind-finite group} instead of ind-\(p\)-group.
\end{statement}

One verifies immediately:

\begin{statement}{Proposition 1.4}
\begin{enumerate}[label=(\roman*)]
\item A subgroup and a quotient of an ind-\(p\)-group is again an ind-\(p\)-group.
\item A filtered inductive limit of ind-\(p\)-groups is an ind-\(p\)-group.
\item A finite projective limit of ind-\(p\)-groups is an ind-\(p\)-group.
\end{enumerate}
\end{statement}

\begin{statement}{Definition 1.5}
Let \(X\) be a scheme. A sheaf \(F\) of groups on \(X\) is called a \emph{sheaf of ind-\(p\)-groups}, respectively a sheaf of \emph{ind-finite groups}, if for every \'etale \(U\to X\), with \(U\) quasi-compact, the group \(F(U)\) is an ind-\(p\)-group, respectively an ind-finite group.
\end{statement}

\begin{statement}{Proposition 1.6}
Let \(F\) be a sheaf of groups on a scheme \(X\).
\begin{enumerate}[label=(\roman*)]
\item The sheaf \(F\) is a sheaf of ind-\(p\)-groups if and only if for every geometric point \(\xi\) of \(X\), the fibre \(F_\xi\) is an ind-\(p\)-group.
\item If \(f:X\to X'\) is a morphism and \(F'\) is a sheaf of ind-\(p\)-groups on \(X'\), then \(f^*F'\) is a sheaf of ind-\(p\)-groups.
\item If \(f:X\to X'\) is quasi-compact and \(F\) is a sheaf of ind-\(p\)-groups on \(X\), then \(f_*F\) is a sheaf of ind-\(p\)-groups.
\end{enumerate}
\end{statement}

\begin{prooftext}
For (i), suppose first that \(F\) is a sheaf of ind-\(p\)-groups, and let \(\xi\) be a geometric point of \(X\). Then \(F_\xi=\varinjlim F(X')\), where \(X'\) runs through a pseudo-filtered system of schemes \'etale over \(X\) (VIII, 3.9). The affine, hence quasi-compact, \(X'\) form a cofinal system, so \(F_\xi\) is an inductive limit of ind-\(p\)-groups, hence an ind-\(p\)-group by Proposition 1.4(ii). Conversely, assume that every fibre \(F_\xi\) is an ind-\(p\)-group, and let \(U\to X\) be \'etale with \(U\) quasi-compact. Let \(S\subset F(U)\) be finite. For every geometric point \(\xi\) of \(U\), the image of \(S\) in \(F_\xi\) generates a finite group. Since a finite group is finitely presented, one easily concludes that there exists an \'etale \(U_\xi\) over \(U\), pointed by \(\xi\), such that \(S\) generates a finite group in \(F(U_\xi)\) (\cf VIII, 4). A finite number of such \(U_\xi\) cover \(U\), say \(\{U_1,\ldots,U_n\}\). Then the image of \(S\) in \(\prod F(U_i)\) generates a finite subgroup, and since \(F(U)\subset\prod F(U_i)\), the conclusion follows. Assertion (ii) is trivial from (i), and (iii) is immediate.
\end{prooftext}
